\subsubsection{2.2.5 Emergenz der Raumzeit durch Projektion}

Nachdem die numerischen Untersuchungen zunehmend darauf hindeuteten, dass der OPP selbst keine intrinsische vierdimensionale Struktur besitzt, rückte eine neue Fragestellung in den Mittelpunkt:

\begin{center}
\emph{Entsteht die Raumzeit erst durch den Beobachtungsprozess beziehungsweise durch die Projektion?}
\end{center}

Dies stellt einen grundlegenden Perspektivwechsel gegenüber klassischen physikalischen Theorien dar.

In der Allgemeinen Relativitätstheorie existiert die Raumzeit als fundamentales Objekt. Materie und Energie verändern ihre Geometrie, doch die Raumzeit selbst gehört bereits zur ontologischen Grundausstattung der Theorie.

Die PST verfolgt dagegen eine radikalere Hypothese:

\[
\boxed{
\text{Raumzeit ist nicht fundamental, sondern emergent.}
}
\]

Der OPP besitzt zunächst keinerlei ausgezeichnete Raumrichtung, keine Metrik und keine klassische Zeitkoordinate.

Er bildet lediglich ein hochdimensionales Netz von Relationen.

\section*{Der OPP besitzt keine Raumzeit}

Die numerischen Untersuchungen lieferten hierfür mehrere Hinweise.

Der rohe OPP zeigte effektive Dimensionen zwischen ungefähr

\[
30
\le
d_{\mathrm{eff}}
\le
45.
\]

Diese Dimension blieb gegenüber Änderungen der Punktzahl bemerkenswert stabil.

Insbesondere konnte keine natürliche vierdimensionale Struktur gefunden werden.

Dies spricht dafür, dass

\[
\boxed{
\text{die Raumzeit keine Eigenschaft des OPP selbst ist.}
}
\]

Vielmehr scheint sie erst während der Projektion zu entstehen.

\section*{Die Projektion als Informationsfilter}

Der Begriff der Projektion erhielt dadurch eine wesentlich tiefere Bedeutung.

Bisher wurde sie als mathematische Dimensionsreduktion verstanden.

Die Diskussion führte jedoch zu einer allgemeineren Interpretation.

Eine Projektion ist kein bloßer Rechenalgorithmus.

Sie stellt vielmehr einen Informationsfilter dar.

Formal

\[
\Pi :
\mathcal{O}
\rightarrow
\mathcal{M},
\]

wobei

\begin{itemize}

\item
\(
\mathcal{O}
\)
den vollständigen OPP beschreibt,

\item
\(
\mathcal{M}
\)
die beobachtbare Welt.

\end{itemize}

Dabei geht Information verloren.

Gerade dieser Informationsverlust erzeugt möglicherweise die klassische Welt.

\section*{Beobachterabhängigkeit}

Ein weiterer Gedanke ergab sich aus der Analogie zur Wahrnehmung biologischer Systeme.

Der Mensch besitzt nur einen kleinen Ausschnitt aller möglichen Sinnesinformationen.

Beispielsweise

\begin{itemize}

\item hört er nur einen kleinen Frequenzbereich,

\item sieht nur einen winzigen Ausschnitt des elektromagnetischen Spektrums,

\item besitzt eine endliche räumliche und zeitliche Auflösung.

\end{itemize}

Dennoch erscheint ihm daraus eine konsistente Welt.

Dies führte zu einer Analogie.

Vielleicht verhält sich die Raumzeit genauso.

Nicht der OPP besitzt notwendigerweise vier Dimensionen,

sondern

\[
\boxed{
\text{der Beobachter besitzt nur Zugriff auf eine vierdimensionale Projektion.}
}
\]

Die Raumzeit wäre damit keine Eigenschaft des Universums an sich,

sondern eine Eigenschaft des Beobachtungsprozesses.

\section*{Projektion und Quantisierung}

Bereits früh wurde im OP die Vermutung formuliert,

dass Quantisierung nicht notwendigerweise fundamental sein muss.

Sie könnte aus begrenzter Informationsauflösung entstehen.

Die Projektion liefert hierfür einen gemeinsamen Rahmen.

Endliche Informationsauflösung erzeugt

\begin{itemize}

\item diskrete Messwerte,

\item scheinbare Quantisierung,

\item klassische Objektgrenzen.

\end{itemize}

Die gleiche Projektion könnte gleichzeitig auch die Raumzeit erzeugen.

Damit würden zwei bislang getrennte Probleme auf dieselbe Ursache zurückgeführt.

\[
\boxed{
\text{Dieselbe Projektion erzeugt sowohl Quantisierung als auch Raumzeit.}
}
\]

Diese Idee besitzt innerhalb der PST besondere Attraktivität, da sie zwei bislang unabhängige Annahmen auf ein einziges Prinzip reduziert.

\section*{Lorentz-Invarianz als Eigenschaft der Projektion}

Eine der größten Herausforderungen jeder emergenten Raumzeittheorie besteht darin,

die Lorentz-Invarianz zu erklären.

Falls Raumzeit lediglich eine Projektion darstellt,

muss diese Projektion automatisch Lorentz-invariant sein.

Die Konsequenz lautet:

Nicht der OPP besitzt notwendigerweise eine Lorentz-Metrik,

sondern

\[
\boxed{
\Pi
\text{ erzeugt eine Lorentz-invariante Beschreibung.}
}
\]

Dadurch würde verständlich,

warum sämtliche Beobachter dieselben physikalischen Gesetze messen,

obwohl der zugrunde liegende OPP wesentlich allgemeiner aufgebaut ist.

\section*{Die Analogie zur Computergrafik}

Zur Veranschaulichung wurde eine Analogie zur Computergrafik verwendet.

Ein dreidimensionales Objekt besitzt unendlich viele mögliche Projektionen auf eine Ebene.

Je nach Blickwinkel entstehen völlig unterschiedliche Bilder.

Das Objekt selbst bleibt unverändert.

Übertragen auf die PST ergibt sich

\[
\text{OPP}
\quad
\leftrightarrow
\quad
3D\text{-Objekt},
\]

während

\[
\text{Raumzeit}
\quad
\leftrightarrow
\quad
2D\text{-Projektion}.
\]

Nur handelt es sich im OPP nicht um drei,

sondern um viele Dutzend Dimensionen.

Die beobachtete Physik wäre lediglich eine besonders ausgezeichnete Projektion.

\section*{Die Auswahl der Projektion}

Die numerischen Ergebnisse zeigten,

dass verschiedene Projektionsverfahren sehr unterschiedliche Resultate liefern.

Unter anderem wurden untersucht

\begin{itemize}

\item PCA,

\item Kernel-PCA,

\item Random Projection,

\item Diffusion Maps.

\end{itemize}

Besonders Diffusion Maps und Kernel-PCA lieferten stabile effektive Dimensionen nahe vier.

Dies führte zu der Vermutung,

dass nicht jede mathematische Projektion physikalisch sinnvoll ist.

Vielmehr existiert offenbar eine Klasse bevorzugter Projektionen.

Die zentrale Forschungsfrage lautet daher

\[
\boxed{
\text{Welche Eigenschaften zeichnen physikalisch zulässige Projektionen aus?}
}
\]

\section*{Eine neue Interpretation der Emergenz}

Die Diskussion führte schließlich zu einer allgemeineren Sichtweise.

Emergenz bedeutet in der PST nicht,

dass aus einfachen Bausteinen komplizierte Strukturen aufgebaut werden.

Stattdessen bedeutet Emergenz,

dass

\[
\boxed{
\text{durch Informationsreduktion neue physikalische Begriffe entstehen.}
}
\]

Raum,

Zeit,

Objekte,

Teilchen,

Felder,

ja möglicherweise sogar Kausalität,

wären keine fundamentalen Bestandteile der Wirklichkeit,

sondern stabile Begriffe,

die erst nach der Projektion auftreten.

\section*{Ausblick}

Die Projektion entwickelt sich damit zum eigentlichen Bindeglied zwischen Ontologie und Physik.

Der OPP beschreibt,

\emph{was existiert}.

Die Projektion beschreibt,

\emph{was davon beobachtbar wird}.

Die beobachtete Physik entsteht schließlich als stabile Struktur dieser Projektion.

Formal lässt sich dieser Zusammenhang zusammenfassen als

\[
\boxed{
\mathcal{O}
\;
\overset{\Pi}{\longrightarrow}
\;
(\text{Raumzeit},
\text{Materie},
\text{Wechselwirkungen},
\text{Quantisierung}).
}
\]

Damit wird die Suche nach geeigneten Projektionsoperatoren zu einer der zentralen Aufgaben der weiteren Entwicklung der PST.

